\chapter{Graphs: counting, paths, and extremal choices}\label{ch:graphs}
\goal{Describe a pattern of connections as a graph. Prove an identity by counting one collection of objects in two different ways. Use the pigeonhole principle and averages to show that something must exist. Choose an object that is as short or as long as possible, and show that this choice has the property we want. Learn the language of graphs that the chapters on matchings build on, and the language of directed graphs needed to read the exact statement of a research problem.}

\section{A graph is a pattern of connections}\label{ch07:sec:graph}
Many situations consist of objects, some pairs of which are connected: people who have shaken hands, towns joined by roads. A graph keeps only which objects there are and which pairs of them are connected, and sets aside everything else, such as how long a road is.

\par\noindent\begin{minipage}{\linewidth}
\begin{definition}[Graph]\label{ch07:def:graph}
A \emph{graph} consists of a finite set $V$, whose members are called \emph{vertices}, and a set $E$, whose members are called \emph{edges}. Each edge is a set $\{u,v\}$ of two different vertices $u$ and $v$. These two vertices are the \emph{endpoints} of the edge, and we say that the edge \emph{joins} them.
\end{definition}
\end{minipage}\par

The full name of this kind of graph is a \emph{finite simple undirected graph}, and each word records a choice.
\begin{itemize}
\item \emph{Finite}: the set $V$ is finite (Section~\ref{ch03:sec:sets}), and therefore so is $E$, because a finite set has only finitely many pairs of members.
\item \emph{Undirected}: since $\{u,v\}=\{v,u\}$, an edge does not run from one endpoint to the other; it simply joins them.
\item \emph{Simple}: there are no \emph{loops}, edges joining a vertex to itself, because the two endpoints of an edge are different. And two vertices are joined by at most one edge, because $E$ is a set: listing $\{u,v\}$ twice does not give it a second member, just as $\{1,1,3\}=\{1,3\}$.
\end{itemize}
In this book, the word \emph{graph} on its own always means a graph in the sense of Definition~\ref{ch07:def:graph}; graphs whose edges have a direction are called \emph{directed} graphs and appear in Section~\ref{ch07:sec:sink}.

We usually write the edge $\{u,v\}$ as $uv$, or equally as $vu$. (In Section~\ref{ch01:sec:letters} two letters side by side meant a product; here they mean an edge, and the context always makes clear which is intended.)

Here is an example that we will use throughout the chapter. Its vertex set is $V=\{a,b,c,d\}$, and its edge set is
\[
E=\{ab,\ bc,\ ca,\ cd\}.
\]
We call it the \emph{example graph} $G$. In a drawing, each vertex is a small circle, and each edge is a line between its two endpoints:
\begin{center}
\begin{tikzpicture}[every node/.style={circle,draw=ink,minimum size=6mm,inner sep=1pt},line width=.7pt]
\node(a)at(0,0){$a$};\node(b)at(2,0){$b$};\node(c)at(1,1.5){$c$};\node(d)at(3,1.5){$d$};
\draw(a)--(b)--(c)--(a);\draw(c)--(d);
\end{tikzpicture}
\end{center}
A drawing is only a picture of a graph: the positions of the circles, the shapes of the lines, and any points where two lines cross carry no information. In particular, a crossing point is not a vertex, and moving the circle for $d$ to the far left would give another picture of the same graph.

Two vertices $u$ and $v$ are \emph{neighbors}, or \emph{adjacent}, when $uv$ is an edge. The \emph{degree} of a vertex $v$, written $\deg(v)$, is the number of its neighbors. In the example graph, the neighbors of $c$ are $a$, $b$, and $d$, so $\deg(c)=3$. The four degrees are
\[
\deg(a)=2,\qquad\deg(b)=2,\qquad\deg(c)=3,\qquad\deg(d)=1.
\]
A vertex with no neighbors has degree $0$; it still belongs to the graph.

The degree of $v$ is also the number of edges at $v$, that is, edges with $v$ as an endpoint: each neighbor $w$ gives one such edge, $vw$, different neighbors give different edges, and every edge at $v$ arises in this way from its other endpoint. For example, the edges at $c$ are $ca$, $cb$, and $cd$. This uses simplicity: if two edges could join the same pair of vertices, one neighbor could account for two edges.

\pauseq{In a graph with $n$ vertices, can a vertex have degree larger than $n-1$?}{No. The neighbors of $v$ are different from $v$, because the two endpoints of an edge are different, so they lie among the other $n-1$ vertices, and $\deg(v)\le n-1$. Equality holds exactly when $v$ is joined to every other vertex, as $c$ is in the example graph.}

\section{Double counting and the handshake theorem}\label{ch07:sec:double}
Add up the degrees in the example graph $G$, which has four edges: $2+2+3+1=8=2\cdot4$. This is no accident.

Consider every pair $(v,e)$ in which $v$ is a vertex, $e$ is an edge, and $v$ is an endpoint of $e$, that is, $v\in e$. The table below has one row for each vertex of $G$ and one column for each edge; an entry $1$ marks a pair $(v,e)$ with $v\in e$.
\begin{center}
\begin{tabular}{@{}lccccc@{}}
\toprule
 & $ab$ & $bc$ & $ca$ & $cd$ & row total\\
\midrule
$a$ & $1$ & $0$ & $1$ & $0$ & $2$\\
$b$ & $1$ & $1$ & $0$ & $0$ & $2$\\
$c$ & $0$ & $1$ & $1$ & $1$ & $3$\\
$d$ & $0$ & $0$ & $0$ & $1$ & $1$\\
\midrule
column total & $2$ & $2$ & $2$ & $2$ & $8$\\
\bottomrule
\end{tabular}
\end{center}
Count the entries equal to $1$ in two ways. The row of a vertex has a $1$ for each edge at that vertex, so the row totals are the degrees, which add up to $8$. Every edge has two endpoints, so every column total is $2$, and the columns also add up to $8$: both totals count the same entries.

Counting one collection of objects in two ways and setting the two answers equal is called \emph{double counting}. Its essential step is to describe the collection precisely, so that we can check that each grouping counts every member exactly once.

To state the general result, we extend the sum notation of Section~\ref{ch01:sec:notation}: $\sum_{v\in V}\deg(v)$ means that we add $\deg(v)$ once for each vertex $v$ in $V$; for the example graph it is $\deg(a)+\deg(b)+\deg(c)+\deg(d)=8$. The order of the terms does not matter, because the rules $a+b=b+a$ and $(a+b)+c=a+(b+c)$ of Section~\ref{ch01:sec:letters} let us reorder and regroup the terms of any finite sum. In the same way, $\sum_{e\in E}2$ adds $2$ once for each edge, so it equals $2|E|$.

\par\noindent\begin{minipage}{\linewidth}
\begin{theorem}[Handshake theorem]\label{thm:handshake}
In every graph with vertex set $V$ and edge set $E$,
\[
\sum_{v\in V}\deg(v)=2|E|.
\]
\end{theorem}
\end{minipage}\par

The name comes from a party: if the guests are vertices and the handshakes are edges, then $\deg(v)$ is the number of hands that $v$ shook, and each handshake is counted twice in the total, once by each of the two people involved.

\begin{proof}
Let $P$ be the set of all pairs $(v,e)$ with $v\in V$, $e\in E$, and $v\in e$. We count $P$ in two ways.

First group the pairs by their vertex. For a fixed vertex $v$, the pairs in $P$ with first entry $v$ are the pairs $(v,e)$ in which $e$ is an edge at $v$, and by Section~\ref{ch07:sec:graph} there are exactly $\deg(v)$ such edges. Each pair in $P$ has just one first entry, so it lies in exactly one of these groups, and $|P|=\sum_{v\in V}\deg(v)$.

Next group the pairs by their edge. For a fixed edge $e=\{u,w\}$, the pairs in $P$ with second entry $e$ are $(u,e)$ and $(w,e)$, two different pairs because $u\ne w$. Again every pair lies in exactly one group, so $|P|=\sum_{e\in E}2=2|E|$.

Both expressions equal $|P|$, so they are equal to each other.
\end{proof}

In the example graph, two vertices, $c$ and $d$, have odd degree. That this number is even is no coincidence either.

\par\noindent\begin{minipage}{\linewidth}
\begin{proposition}\label{ch07:prop:odd}
In every graph, the number of vertices of odd degree is even.
\end{proposition}
\end{minipage}\par

\begin{proof}
Let $k$ be the number of vertices of odd degree. By Lemma~\ref{ch04:lem:even-or-odd}, every degree is either even or odd, and not both. So each vertex contributes its degree to exactly one of two sums: $S_{\mathrm{even}}$, over the vertices of even degree, and $S_{\mathrm{odd}}$, over the vertices of odd degree. Together they make up the degree sum, so $S_{\mathrm{even}}+S_{\mathrm{odd}}=2|E|$ by the handshake theorem.
\begin{itemize}
\item $S_{\mathrm{even}}$ is even, because a factor $2$ can be taken out of every term. (If no vertex has even degree, $S_{\mathrm{even}}$ is the empty sum $0$, which is even too.)
\item Therefore $S_{\mathrm{odd}}=2|E|-S_{\mathrm{even}}$ is even, since $2x-2y=2(x-y)$.
\end{itemize}
Now write the degrees of the $k$ vertices of odd degree, using Definition~\ref{def:parity}, as $2\ell_1+1,\ldots,2\ell_k+1$ with integers $\ell_1,\ldots,\ell_k$. Adding them gives
\[
S_{\mathrm{odd}}=(2\ell_1+1)+\cdots+(2\ell_k+1)=2(\ell_1+\cdots+\ell_k)+k,
\]
because each of the $k$ terms contributes one $1$. Hence $k=S_{\mathrm{odd}}-2(\ell_1+\cdots+\ell_k)$ is the difference of two even numbers, so $k$ is even. (If $k=0$, the conclusion holds at once, since $0$ is even.)
\end{proof}

For instance, the proposition shows that seven people at a party cannot each shake hands with exactly three of the others: all seven vertices would have the odd degree $3$, and $7$ is odd. (The handshake theorem shows this directly: the degree sum $7\cdot3=21$ would be odd, but it equals $2|E|$.)

\section{The pigeonhole principle and the average}\label{ch07:sec:pigeonhole}
If many objects are placed in a few boxes, some box must be crowded. Here is the precise statement, one of the most useful facts in counting.

\par\noindent\begin{minipage}{\linewidth}
\begin{proposition}[Pigeonhole principle]\label{ch07:prop:pigeonhole}
Let $n$, $r$, and $k$ be natural numbers. If $n$ objects are placed in $r$ boxes and $n>rk$, then some box contains more than $k$ objects.
\end{proposition}
\end{minipage}\par

\begin{proof}
We prove the contrapositive (Section~\ref{ch02:sec:contrapositive}). By Section~\ref{ch02:sec:negation}, the negation of the conclusion is ``every box contains at most $k$ objects,'' and the negation of $n>rk$ is $n\le rk$. So suppose that every box contains at most $k$ objects. Each object lies in exactly one box, so $n$ is the sum of the numbers of objects in the $r$ boxes. Each of these $r$ numbers is at most $k$, so their sum is at most $rk$, that is, $n\le rk$.
\end{proof}

The name pictures pigeons settling into the holes of a pigeon house: if there are more pigeons than holes, some hole holds at least two pigeons. This is the case $k=1$, the one used most often. For example, among any $13$ people, two have their birthdays in the same month: the people are the objects, the $12$ months are the boxes, and $13>12\cdot1$. Larger values of $k$ work in the same way. Among any $25$ people, some month contains the birthdays of at least three of them, because $25>12\cdot2$. In each case the principle tells us that such a month \emph{exists}, but not which month it is.

The pigeonhole principle can also be stated in the language of functions from Chapter~\ref{ch:maps}. Placing each object in a box defines a function from the set of objects to the set of boxes. In that language, the case $k=1$ is the fact from Section~\ref{ch03:sec:injsurj} that a function from a finite set $A$ to a finite set $B$ with $|A|>|B|$ cannot be injective: two different inputs must share an output.

\pauseq{A class has $30$ students. Must some month contain the birthdays of at least three of them? Of at least four?}{Take the $12$ months as boxes. Since $30>12\cdot2=24$, the pigeonhole principle with $k=2$ shows that some month contains more than two birthdays, that is, at least three. For four birthdays we would need $k=3$, but $30$ is not larger than $12\cdot3=36$, so the principle gives no guarantee. In fact no guarantee is possible: six months with three birthdays each and six months with two each account for $6\cdot3+6\cdot2=30$ students, and no month holds four.}

A closely related fact concerns averages. The \emph{average} of real numbers $a_1,\ldots,a_n$, where $n\ge1$, is
\[
\bar a=\frac{a_1+\cdots+a_n}{n}.
\]
(The bar over the $a$ is part of the name; read $\bar a$ as ``$a$ bar.'')

\par\noindent\begin{minipage}{\linewidth}
\begin{proposition}[Averaging principle]\label{ch07:prop:average}
Let $a_1,\ldots,a_n$ be real numbers, where $n\ge1$, and let $\bar a$ be their average. Then at least one of the numbers is at least $\bar a$, and at least one of them is at most $\bar a$.
\end{proposition}
\end{minipage}\par

\begin{proof}
Suppose, for contradiction, that every $a_i$ is strictly less than $\bar a$. Adding these $n$ strict inequalities gives
\[
a_1+\cdots+a_n<n\bar a.
\]
(Two strict inequalities $x<y$ and $x'<y'$ can be added: $x+x'<y+x'<y+y'$, by the rule for adding the same number to both sides from Section~\ref{ch02:sec:quantifiers}. Repeating this adds all $n$ of them. Since $n\ge1$, there is at least one inequality to add, so the result is strict.) But the definition of the average says exactly that $n\bar a=a_1+\cdots+a_n$, so this number would be strictly less than itself, which is impossible. Hence some $a_i$ satisfies $a_i\ge\bar a$.

The same argument with every inequality reversed proves the second claim: if every $a_i$ were strictly greater than $\bar a$, adding would give $a_1+\cdots+a_n>n\bar a$, again impossible.
\end{proof}

For instance, the numbers $1,3,8$ have average $12/3=4$. The number $8$ is at least $4$, and the numbers $1$ and $3$ are at most $4$. The condition $n\ge1$ is needed for the average to make sense at all: with $n=0$ the formula would divide by zero, and there would be no number to choose.

The pigeonhole principle is a special case of the averaging principle. If $n$ objects lie in $r\ge1$ boxes, the numbers of objects in the boxes add up to $n$, so their average is $n/r$. When $n>rk$, this average is larger than $k$, and the averaging principle gives a box containing at least $n/r$ objects, which is more than $k$.

The words ``at least one'' are essential: an average does not control every number. Apply the averaging principle to the degrees of a graph with $n\ge1$ vertices and $m$ edges. By the handshake theorem the degrees add up to $2m$, so their average is $2m/n$, and therefore some vertex has degree at least $2m/n$. In the example graph $G$, with $n=m=4$, the average degree is $2$, and the vertex $c$ has degree $3\ge2$, as promised. Nothing is promised about the other degrees; Exercise~\ref{ex:7.5} asks how far from the average they can be. Chapter~\ref{ch:probability} develops this observation into a method for proving that an object exists: compute an average over random choices, and conclude that some choice is at least as good as the average.

\section{Walks, paths, and shortest objects}\label{ch07:sec:walks}
A \emph{walk} from a vertex $u$ to a vertex $w$ is a finite sequence of vertices $v_0,v_1,\ldots,v_k$ that starts at $v_0=u$, ends at $v_k=w$, and in which every two consecutive vertices are joined by an edge: $v_{i-1}v_i$ is an edge for every $i=1,\ldots,k$. Its \emph{length} is $k$, the number of steps, one less than the number of entries. A walk may visit a vertex or use an edge more than once; each use of an edge counts as a separate step. A \emph{path} is a walk in which no vertex appears twice. A single vertex $v_0$ counts as a walk and a path of length $0$ from $v_0$ to itself.

In the example graph $G$, the sequence $a,c,b,c,d$ is a walk of length $4$ from $a$ to $d$, since $ac$, $cb$, $bc$, and $cd$ are edges. It visits $c$ twice, so it is not a path. The sequence $a,c,d$ is a path of length $2$ from $a$ to $d$, and $a,d$ is not a walk at all, because $ad$ is not an edge. Since edges have no direction, reading a walk backward gives a walk in the opposite direction: $d,c,a$ is a walk from $d$ to $a$. We call two vertices \emph{connected} when there is a walk from one to the other.

A walk that repeats a vertex contains a detour that can be cut out. Between its two visits to $c$, the walk $a,c,b,c,d$ makes the round trip $c,b,c$; deleting it, but keeping one copy of $c$, leaves the shorter walk $a,c,d$. The step to $d$ now follows the first visit to $c$ instead of the second, which is legitimate because both visits are at the same vertex. The next lemma uses this cut to show that a shortest walk never repeats a vertex.

\par\noindent\begin{minipage}{\linewidth}
\begin{lemma}\label{ch07:lem:shortest-walk}
Let $u$ and $w$ be vertices of a graph, and suppose that there is a walk from $u$ to $w$. Then there is a walk from $u$ to $w$ of smallest possible length, and every such shortest walk is a path. In particular, there is a path from $u$ to $w$.
\end{lemma}
\end{minipage}\par

\begin{proof}
The lengths of walks from $u$ to $w$ form a set of natural numbers, which is nonempty because at least one such walk exists. By the well-ordering property (Section~\ref{ch04:sec:contrapositive}, extended to natural numbers in Section~\ref{ch05:sec:not-circular}), this set has a smallest member $L$. Let $v_0,v_1,\ldots,v_L$ be any walk from $u$ to $w$ of this length.

Suppose, for contradiction, that $v_i=v_j$ for some positions $i<j$. Delete the entries $v_{i+1},\ldots,v_j$, leaving the sequence
\[
v_0,\ \ldots,\ v_i,\ v_{j+1},\ \ldots,\ v_L
\]
(if $j=L$, the sequence simply ends at $v_i$). This is a walk from $u$ to $w$:
\begin{itemize}
\item Each pair of consecutive entries is joined by an edge. Every such pair was already consecutive in the original walk, except for the new pair $v_i,v_{j+1}$, which occurs only when $j<L$. Since $v_i=v_j$, the edge $v_iv_{j+1}$ is the same as $v_jv_{j+1}$, which is an edge of the original walk.
\item It starts at $v_0=u$.
\item It ends at $w$: if $j<L$, its last entry is $v_L=w$, and if $j=L$, its last entry is $v_i=v_j=v_L=w$.
\end{itemize}
Its length is $L-(j-i)$, because $j-i$ entries were deleted. Since $j>i$, this is less than $L$, which contradicts the choice of $L$. Hence no vertex appears twice. Since the shortest walk was arbitrary, every shortest walk from $u$ to $w$ is a path, and in particular such a path exists.
\end{proof}

When $u$ and $w$ are connected, the \emph{distance} from $u$ to $w$ is the length of a shortest walk from $u$ to $w$. By Lemma~\ref{ch07:lem:shortest-walk}, it is also the length of a shortest path, because a shortest walk is a path and every path is a walk. In the example graph, the distance from $a$ to $d$ is $2$: the path $a,c,d$ has length $2$, and no shorter walk exists, because a walk of length $0$ would need $a=d$ and a walk of length $1$ would need $ad$ to be an edge. Establishing a distance takes both facts: a walk of that length, and a reason why no shorter walk exists.

The proof of Lemma~\ref{ch07:lem:shortest-walk} never built a path step by step. It chose a walk that was as short as possible and showed that this choice already has the property we wanted. An object chosen to be as small or as large as possible is called \emph{extremal}, and an argument that starts from such a choice is called an \emph{extremal argument}. We have met this idea before: Section~\ref{ch04:sec:sqrt2} chose a fraction for $\sqrt2$ with the smallest possible positive denominator, and Section~\ref{ch05:sec:not-circular} chose a smallest counterexample.

\section{Directed edges and a guaranteed sink}\label{ch07:sec:sink}
Some connections have a direction: a one-way street, a link from one web page to another, a game won by one player against another. We record such a connection as an arrow.

\par\noindent\begin{minipage}{\linewidth}
\begin{definition}[Directed and oriented graphs]\label{ch07:def:directed}
A \emph{directed graph} consists of a set $V$ of vertices and a set of \emph{directed edges}. Each directed edge is an ordered pair $(u,v)$ of two different vertices; we write it $u\to v$ and draw it as an arrow from $u$ to $v$. An \emph{oriented graph} is a directed graph in which no two vertices are joined in both directions: $u\to v$ and $v\to u$ never both occur.
\end{definition}
\end{minipage}\par

Because the entries of an ordered pair come in a definite order (Section~\ref{ch03:sec:sets}), $u\to v$ and $v\to u$ are different directed edges. As in a graph, there are no loops $u\to u$. Some pairs of vertices may have no edge in either direction. Unlike a graph, a directed graph may have infinitely many vertices, so we say ``finite'' explicitly whenever a result needs it.

We also use $u\to w$ as a statement, saying that $(u,w)$ is a directed edge; $u\to w$ \emph{fails} when it is not. (This arrow joins two vertices; it is not the arrow of a function $f:A\to B$.) If $u\to w$, we call $w$ an \emph{out-neighbor} of $u$. The \emph{outdegree} of $u$ is the number of its out-neighbors, that is, the number of arrows leaving $u$. A \emph{sink} is a vertex with outdegree $0$: no arrow leaves it.

Walks and paths now have to follow the arrows. A \emph{directed walk} from $v_0$ to $v_k$ is a sequence of vertices $v_0,v_1,\ldots,v_k$ with $v_{i-1}\to v_i$ for every $i=1,\ldots,k$; its length is $k$. A \emph{directed path} is a directed walk with no repeated vertex. A \emph{directed cycle} is a directed walk $v_0,v_1,\ldots,v_k$ with at least one step, $k\ge1$, that returns to its start, $v_k=v_0$, and has no other repetition: the vertices $v_0,v_1,\ldots,v_{k-1}$ are all different. A directed graph is \emph{acyclic} when it has no directed cycle.

In an oriented graph, a directed cycle passes through at least three vertices, because a cycle with one step would be a loop and a cycle with two steps would need arrows in both directions. Here are two oriented graphs, each with vertices $a$, $b$, $c$:
\begin{center}
\begin{tikzpicture}[>=Stealth,every node/.style={circle,draw=ink,minimum size=6mm,inner sep=1pt},line width=.7pt]
\node(a)at(0,0){$a$};\node(b)at(1.7,0){$b$};\node(c)at(3.4,0){$c$};
\draw[->](a)--(b);\draw[->](b)--(c);
\node(a2)at(5.6,0){$a$};\node(b2)at(7.3,0){$b$};\node(c2)at(6.45,1.4){$c$};
\draw[->](a2)--(b2);\draw[->](b2)--(c2);\draw[->](c2)--(a2);
\end{tikzpicture}
\end{center}
On the left is the \emph{chain} $a\to b\to c$. It is acyclic, since every arrow leads further along the chain, and its only sink is $c$. On the right is the directed cycle $a\to b\to c\to a$. Replacing its arrow $c\to a$ by $a\to c$ gives a triangle with no directed cycle, since then every arrow goes to a later letter.

Does every directed graph have a sink? A directed graph with no vertices has none. An infinite directed graph can fail too: take the natural numbers as vertices, with an edge $n\to n+1$ for every $n$. There is no sink, yet the graph is acyclic, because the numbers strictly increase along any directed walk. Exercise~\ref{ex:7.3} asks what can happen when the graph is finite but not acyclic.

\par\noindent\begin{minipage}{\linewidth}
\begin{theorem}\label{thm:sink}
Every nonempty finite acyclic oriented graph has a sink.
\end{theorem}
\end{minipage}\par

Here \emph{nonempty} means that there is at least one vertex. The theorem promises \emph{one} sink, not that every vertex is a sink; in the chain, only $c$ is. The word ``oriented'' excludes no acyclic directed graph: a directed graph with both $u\to v$ and $v\to u$ has the directed cycle $u,v,u$. We keep the word because the conjecture in Section~\ref{ch07:sec:seymour} is about oriented graphs.

\noindent\textit{Plan.} Take a directed path that is as long as possible, and look at its last vertex. An arrow leaving that vertex would either reach a new vertex, which would give a longer path, or return to the path, which would close a directed cycle.
\begin{proof}
A directed path never repeats a vertex, so it has at most $|V|$ vertices and therefore length at most $|V|-1$. (By the pigeonhole principle, a list of more than $|V|$ vertices contains some vertex twice.) Since the graph is nonempty, it has a vertex, and a single vertex is a directed path of length $0$. So the lengths of the directed paths form a nonempty subset of $\{0,1,\ldots,|V|-1\}$, which has a largest member. Choose a directed path of this greatest length, and write it $v_0\to v_1\to\cdots\to v_k$.

Suppose, for contradiction, that its last vertex $v_k$ has an out-neighbor $w$.
\begin{itemize}
\item If $w$ is not on the path, then $v_0\to v_1\to\cdots\to v_k\to w$ is a directed path of length $k+1$, which contradicts the choice of a path of greatest length.
\item If $w$ is on the path, say $w=v_j$, then $j\ne k$, because $v_k\to v_k$ would be a loop. So $j<k$. Then $v_j\to v_{j+1}\to\cdots\to v_k\to v_j$ is a directed walk with $k-j+1$ steps that returns to its start, and its vertices $v_j,\ldots,v_k$ before the return are all different. It is a directed cycle, which contradicts acyclicity.
\end{itemize}
Both cases are impossible. Hence $v_k$ has no out-neighbor, so it is a sink.
\end{proof}

Each hypothesis has a visible job in this proof. Nonemptiness supplies a path to choose from. Finiteness guarantees a directed path of greatest length; in the directed graph on the natural numbers, every directed path can be extended. Acyclicity rules out an arrow from the last vertex back into the path.

\section{Tournaments and a vertex of maximum outdegree}\label{ch07:sec:tournament}
In a round-robin competition, every pair of players meets exactly once, and every game has a winner. We can record the results as an oriented graph: the players are the vertices, and $u\to v$ means that $u$ beat $v$.

\par\noindent\begin{minipage}{\linewidth}
\begin{definition}[Tournament]\label{ch07:def:tournament}
A \emph{tournament} is an oriented graph in which every pair of different vertices is joined by exactly one directed edge: for all vertices $u\ne v$, exactly one of $u\to v$ and $v\to u$ holds.
\end{definition}
\end{minipage}\par

Here is a tournament with four players $p$, $q$, $r$, $s$ and the six results $p\to q$, $p\to r$, $q\to r$, $q\to s$, $r\to s$, and $s\to p$, one for each pair:
\begin{center}
\begin{tikzpicture}[>=Stealth,every node/.style={circle,draw=ink,minimum size=6mm,inner sep=1pt},line width=.7pt]
\node(p)at(0,1.8){$p$};\node(q)at(2.4,1.8){$q$};\node(r)at(2.4,0){$r$};\node(s)at(0,0){$s$};
\draw[->](p)--(q);\draw[->](p)--(r);\draw[->](q)--(r);\draw[->](q)--(s);\draw[->](r)--(s);\draw[->](s)--(p);
\end{tikzpicture}
\end{center}
The outdegrees, the numbers of games won, are $2$ for $p$ and $q$, and $1$ for $r$ and $s$.

Call a vertex $v$ a \emph{king} if it can reach every other vertex by a directed path of length at most two: for every other player $u$, either $v$ beat $u$, or $v$ beat someone who beat $u$. In the example, $p$ is a king: it beat $q$ and $r$, and it reaches $s$ through $p\to q\to s$. So is $q$, which beat $r$ and $s$ and reaches $p$ through $q\to s\to p$. The player $r$ is not a king: its only win was against $s$, and $s$ beat only $p$, so within two steps $r$ does not reach $q$. Must every nonempty finite tournament have a king? An extremal choice answers this question.

\par\noindent\begin{minipage}{\linewidth}
\begin{proposition}\label{thm:king}
In a nonempty finite tournament, every vertex of maximum outdegree is a king: it can reach every other vertex by a directed path of length at most two.
\end{proposition}
\end{minipage}\par

A vertex of maximum outdegree always exists, because the outdegrees form a nonempty finite set of natural numbers, which has a largest member. So the proposition shows that every nonempty finite tournament has a king.

\begin{proof}
Let $v$ be any vertex of maximum outdegree, and let $A$ be the set of its out-neighbors, so that the outdegree of $v$ is $|A|$. Let $u$ be any vertex other than $v$. There are two ways in which $u$ might be reached quickly.
\begin{itemize}
\item If $v\to u$, then $u$ is reached in one step.
\item If some $a\in A$ satisfies $a\to u$, then $u$ is reached in two steps by $v\to a\to u$. This is a directed path, because its three vertices are different: $a\ne v$ because $v\to a$ and there are no loops, $u\ne v$ by the choice of $u$, and $a\ne u$ because $a\to u$.
\end{itemize}
Suppose, for contradiction, that neither happens. Since $v\to u$ fails and $u\ne v$, the tournament condition gives $u\to v$. Also $u\notin A$, because $v\to u$ fails, so $u$ is different from every $a\in A$. For each $a\in A$, the tournament condition gives $a\to u$ or $u\to a$, and $a\to u$ has been excluded, so $u\to a$. Thus $u$ beat every member of $A$, and it also beat $v$, which is not in $A$ because a vertex is never its own out-neighbor. So $u$ has at least $|A|+1$ out-neighbors, while $v$ has exactly $|A|$. This contradicts the choice of $v$ as a vertex of maximum outdegree. Therefore one of the two ways always occurs, and $v$ is a king.
\end{proof}

\pauseq{In the example tournament, $s$ won only one game, so it does not have maximum outdegree. Is $s$ a king? What does your answer say about the converse of Proposition~\ref{thm:king}?}{Yes: $s\to p$, $s\to p\to q$, and $s\to p\to r$ reach all three other players within two steps. So the converse, ``every king has maximum outdegree,'' is false, with $s$ as a counterexample: the proposition gives a sufficient condition for being a king, not a necessary one. The kings of the example are $p$, $q$, and $s$.}

\section{Second out-neighborhoods and a research problem}\label{ch07:sec:seymour}
This section states a problem about directed graphs that resisted every attempt at a proof for more than thirty years. Stating it exactly needs one more idea: distance measured along the arrows.

The deletion argument of Lemma~\ref{ch07:lem:shortest-walk} works just as well for directed walks: cutting out the part between two visits to the same vertex leaves a shorter directed walk with the same start and end, because the arrow that used to leave the second visit now leaves the first, which is the same vertex. So, as in the lemma, whenever there is a directed walk from $v$ to $w$, there is a shortest one, and it is a directed path. Its length is called the \emph{directed distance} from $v$ to $w$. Unlike distance in a graph, directed distance depends on the order of the two vertices. In the chain $a\to b\to c$, the directed distance from $a$ to $c$ is $2$, but there is no directed walk from $c$ back to $a$ at all.

Now let $v$ be a vertex of a directed graph with vertex set $V$. Write $N_1^+(v)$ for the set of out-neighbors of $v$, and call it the \emph{first out-neighborhood} of $v$. Its members are exactly the vertices at directed distance $1$ from $v$. The \emph{second out-neighborhood} of $v$ is the set
\[
N_2^+(v)=\{w\in V:\ w\ne v,\ w\notin N_1^+(v),\ \text{and }v\to u\to w\text{ for some vertex }u\}.
\]
The superscript $+$ is a reminder that we follow the arrows outward from $v$, and the subscript counts the steps. In words, $w$ belongs to $N_2^+(v)$ when some two-step walk $v\to u\to w$ reaches it but no shorter walk does: $w$ is not $v$ itself, at distance $0$, and not a first out-neighbor, at distance $1$. So $N_2^+(v)$ is the set of vertices at directed distance exactly two from $v$. It is a set of vertices, not of routes: if several intermediate vertices $u$ lead to the same $w$, then $w$ is still a single member of $N_2^+(v)$ and is counted once in $|N_2^+(v)|$.

In an oriented graph, the condition $w\ne v$ holds automatically, because a route $v\to u\to v$ would need edges in both directions between $v$ and $u$. In a directed graph that is not oriented, such a route can occur, and the condition $w\ne v$ is what prevents $v$ from counting as its own second out-neighbor.

Two small examples show the definition at work. In the oriented graph with edges $a\to b$, $b\to c$, $a\to c$, and $c\to d$, the two-step routes from $a$ are $a\to b\to c$ and $a\to c\to d$. The vertex $c$ is at the end of the first route, but it is not in $N_2^+(a)$, because the edge $a\to c$ already makes it a first out-neighbor. So
\[
N_1^+(a)=\{b,c\},\qquad N_2^+(a)=\{d\},\qquad |N_1^+(a)|=2,\qquad |N_2^+(a)|=1.
\]
In the directed cycle $a\to b\to c\to a$, each vertex has exactly one first and one second out-neighbor; for instance, $N_1^+(a)=\{b\}$ and $N_2^+(a)=\{c\}$. Here nothing is excluded: the only two-step route from $a$ is $a\to b\to c$, and $c$ is neither $a$ nor an out-neighbor of $a$.

In 1990 Paul Seymour proposed the following statement, now known as \emph{Seymour's second-neighborhood conjecture}:
\begin{quote}
Every nonempty finite oriented graph has a vertex $v$ with $|N_2^+(v)|\ge|N_1^+(v)|$.
\end{quote}
In words: some vertex reaches at least as many new vertices in its second step as in its first. In the four-vertex example above, the vertex $a$ fails the inequality, since $1<2$. The vertex $b$ satisfies it, since $N_1^+(b)=\{c\}$ and $N_2^+(b)=\{d\}$, so $1\ge1$. So does $d$, for a simpler reason. It is a sink, and a sink always satisfies the inequality: it has no out-neighbors, so it has no two-step routes either, and both of its out-neighborhoods are empty. So the interesting graphs are those in which every vertex has at least one out-neighbor.

The word ``oriented'' cannot be dropped from the conjecture. Consider the directed graph with two vertices $u$ and $v$ and both edges $u\to v$ and $v\to u$. The only two-step route from $u$ is $u\to v\to u$, which returns to $u$ itself, so $N_1^+(u)=\{v\}$ and $N_2^+(u)=\varnothing$. The same holds for $v$, so neither vertex satisfies the inequality.

For decades the conjecture remained open, although special cases were settled; the case of tournaments, for example, was proved by David Fisher in 1996. In September 2026, a \emph{preprint}, that is, a research paper made public before a journal has reviewed it, claimed a proof of the full conjecture. It was released in OpenAI's public mathematics repository together with a formal proof in the Lean proof assistant~\cite{oai173}. The repository's description of the formalized statement uses the same conventions as this section: the second out-neighborhood consists of the vertices at directed distance exactly two, so $v$ and its first out-neighbors are excluded, and graphs with sinks are allowed. The proof is long and draws on ideas well beyond this course. Section~\ref{ch16:sec:second} shows how the definition of $N_2^+(v)$ is written in Lean, and Section~\ref{ch17:sec:research} describes how to begin reading a formalized proof of this kind.

If a news report or an AI assistant tells you that ``Seymour's conjecture has been proved,'' the first thing to check is \emph{which statement} was proved. A proof that counted two-step routes instead of vertices, or that allowed first out-neighbors to count as second out-neighbors, would establish a weaker statement. Each change can only make the second count larger, and a larger second count makes the inequality easier to satisfy. The definitions in this section are what let you ask that question precisely.

We can already prove one special case completely.

\par\noindent\begin{minipage}{\linewidth}
\begin{proposition}\label{ch07:prop:acyclic-seymour}
Every nonempty finite acyclic oriented graph has a vertex $v$ with $|N_2^+(v)|\ge|N_1^+(v)|$.
\end{proposition}
\end{minipage}\par

\begin{proof}
By Theorem~\ref{thm:sink}, the graph has a sink $v$. As we saw above, both out-neighborhoods of a sink are empty. Hence $|N_2^+(v)|=0\ge0=|N_1^+(v)|$.
\end{proof}

This is a complete proof for acyclic graphs, and no proof at all for graphs with directed cycles, which is where all the difficulty lies. In fact, a counterexample to the conjecture would have to be a nonempty finite oriented graph without a sink, and by Theorem~\ref{thm:sink} such a graph must contain a directed cycle. Keeping track of which graphs a proof covers is part of keeping the statement contract of Habit~2 (Section~\ref{sec:magnitude}): the hypothesis ``acyclic'' belongs to the contract of Proposition~\ref{ch07:prop:acyclic-seymour}, and removing it changes the statement into the conjecture itself.

\section*{Exercises}
Exercises~\ref{ex:7.1}--\ref{ex:7.4} are core exercises, and Exercises~\ref{ex:7.5} and~\ref{ex:7.6} are transfer exercises. Complete solutions are in Appendix~\ref{app:solb}, under ``Week~7.''
\begin{exercise}\label{ex:7.1}Let $G$ be the example graph of Section~\ref{ch07:sec:graph}. (a)~Add the edge $bd$ to $G$. List the degrees of the vertices of the new graph, and find their sum and the number of vertices of odd degree. Check that both answers agree with Theorem~\ref{thm:handshake} and Proposition~\ref{ch07:prop:odd}. (b)~Do the same for the graph obtained from $G$ by deleting the edge $cd$ instead.\end{exercise}
\begin{exercise}\label{ex:7.2}Let $x_1,x_2,x_3,x_4,x_5$ be five integers, not necessarily different. Prove that there are two indices $i\ne j$ for which $x_i$ and $x_j$ leave the same remainder on division by four. Then give four integers that leave four different remainders on division by four, to show that five cannot be replaced by four.\end{exercise}
\begin{exercise}\label{ex:7.3}Show that Theorem~\ref{thm:sink} becomes false if the word ``acyclic'' is removed: find a nonempty finite oriented graph with no sink. Use as few vertices as possible, and prove that every nonempty oriented graph with fewer vertices has a sink. Which step of the proof of Theorem~\ref{thm:sink} breaks down for your graph?\end{exercise}
\begin{exercise}\label{ex:7.4}Consider the oriented graph with vertices $a,b,c,d$ and edges $a\to b$, $a\to c$, $b\to d$, and $c\to d$. Find $N_1^+(v)$ and $N_2^+(v)$ for every vertex $v$, and decide which vertices satisfy $|N_2^+(v)|\ge|N_1^+(v)|$. Be careful with $d$, which can be reached from $a$ along two different routes.\end{exercise}
\begin{exercise}\label{ex:7.5}Consider a graph with $n\ge1$ vertices and $m$ edges. Prove that some vertex has degree at most $2m/n$. Then explain, using an example, why this gives no guarantee that all the degrees are close to the average $2m/n$.\end{exercise}
\begin{exercise}\label{ex:7.6}Let $v$ be a king in a tournament, so that every other vertex can be reached from $v$ by a directed path of length at most two. Does it follow that $|N_2^+(v)|\ge|N_1^+(v)|$? A tournament on three vertices is enough to settle the question.\end{exercise}

